\documentclass[a4paper,12pt]{article}
\usepackage[utf8]{inputenc}
\usepackage{amsmath, amssymb}
\usepackage{geometry}
\usepackage{graphicx}
\usepackage[utf8]{inputenc}
\usepackage[T1]{fontenc}
\usepackage[spanish]{babel}
\usepackage{url}
\geometry{margin=1in}

% Datos de la tarea
\title{\textbf{Planeación:} Case study: Estrategias antifrágiles en el contexto de monopolios.}
\author{Eliud Daniel Santillan Aparicio 415069696 \\ Leonardo Vigueras Salazar 310738244}

\date{} % Elimina la fecha


\begin{document}

\maketitle
\section*{Instrucciones}

El presente plan de trabajo se estructura en cinco etapas que permiten abordar de manera progresiva el objetivo general del proyecto: proponer un modelo que formalice las condiciones sistémicas que inducen dinámicas monopólicas, y evaluar si dicho modelo exhibe características antifrágiles.

\subsubsection*{Etapa I: Fundamentación Teórica y Conceptual}

\begin{itemize}
    \item Realizar una revisión bibliográfica sistemática sobre:
    \begin{itemize}
        \item Modelos de formación de monopolios: económicos, históricos y regulatorios.
        \item Sistemas complejos, redes, y modelos biológicos de ecosistemas.
        \item Teoría de la antifragilidad en contextos económicos.
    \end{itemize}
    \item Definir operacionalmente conceptos clave: monopolio, emergencia, voluntad centralizada, condiciones sistémicas, antifragilidad.
    \item Estudiar modelos dinámicos relevantes: ecuaciones de Lotka-Volterra, redes complejas, sistemas autoorganizados.
\end{itemize}

\subsubsection*{Etapa II: Diseño del Modelo Sistémico}

\begin{itemize}
    \item Identificar variables clave: oferta, demanda, regulación, escalabilidad, efectos de red, asimetrías de información, etc.
    \item Seleccionar el marco modelador más adecuado:
    \begin{itemize}
        \item Sistemas dinámicos no lineales.
        \item Modelos de redes dinámicas.
        \item Modelos basados en agentes (ABM).
        \item Modelos híbridos.
    \end{itemize}
    \item Construir formalmente el modelo inicial con especificación de reglas, estados y parámetros.
    \item Simular escenarios hipotéticos con distintas condiciones iniciales y choques externos.
\end{itemize}

\subsubsection*{Etapa III: Validación Empírica y Entrenamiento}

\begin{itemize}
    \item Recopilar casos históricos y conjuntos de datos reales: Standard Oil, East India Company, Amazon, etc.
    \item Ajustar parámetros del modelo utilizando técnicas de calibración: optimización bayesiana, simulación Monte Carlo, aprendizaje automático.
    \item Validar el modelo:
    \begin{itemize}
        \item Comparar resultados simulados con trayectorias históricas.
        \item Evaluar capacidad predictiva y explicativa.
    \end{itemize}
\end{itemize}

\subsubsection*{Etapa IV: Evaluación de Antifragilidad}

\begin{itemize}
    \item Definir métricas de antifragilidad adaptadas al entorno económico.
    \item Someter el modelo a perturbaciones controladas: choques regulatorios, tecnológicos o financieros.
    \item Clasificar los tipos de monopolio según su respuesta: frágiles, robustos, antifrágiles.
\end{itemize}

\subsubsection*{Etapa V: Conclusión y Publicación}

\begin{itemize}
    \item Redactar el documento final del modelo, con estructura académica: introducción, metodología, resultados, discusión y conclusiones.
    \item Analizar las implicaciones políticas, regulatorias y éticas de los resultados.
    \item Preparar materiales de diseminación: artículo académico, resumen ejecutivo, visualización interactiva del modelo.
\end{itemize}




\end{document}